\subsection{Haces localmente libres}
\label{subsection:0.5.4-fr}

\begin{env}[5.4.1]
\label{0.5.4.1-fr}
Sea $X$ un espacio anillado. Se dice que un $\mathcal{O}_X$-Módulo $\mathcal{F}$
es \emph{localmente libre} si, para todo $x\in X$, existe una vecindad
abierta $U$ de $x$ tal que $\mathcal{F}|U$ es isomorfo a un
$(\mathcal{O}_X|U)$-Módulo de la forma $\mathcal{O}_X^{(I)}|U$, donde $I$ puede
depender de $U$. Si para todo $U$, $I$ es finito, se dice que $\mathcal{F}$ es
\emph{de rango finito}; si para todo $U$, $I$ tiene el mismo número finito de elementos
$n$, se dice que $\mathcal{F}$ es \emph{de rango $n$}. Un
$\mathcal{O}_X$-Módulo localmente libre de rango 1 se llama también
\emph{invertible} (cf. (5.4.3)). Si $\mathcal{F}$ es un
$\mathcal{O}_X$-Módulo localmente libre de rango finito, para todo $x\in X$,
$\mathcal{F}_x$ es un $\mathcal{O}_x$-módulo libre de rango finito $n(x)$, y
existe una vecindad $U$ de $x$ tal que $\mathcal{F}|U$ es de rango $n(x)$;
si $X$ es conexo, $n(x)$ es, por tanto, \emph{constante}.

Es claro que todo haz localmente libre es cuasicoherente, y si
$\mathcal{O}_X$ es un haz coherente de anillos, todo
$\mathcal{O}_X$-Módulo localmente libre de rango finito es coherente.

Si $\mathcal{L}$ es localmente libre,
$\mathcal{L}\otimes_{\mathcal{O}_X}\mathcal{F}$ es un funtor \emph{exacto}
en $\mathcal{F}$ en la categoría de los $\mathcal{O}_X$-Módulos.

Tendremos que considerar sobre todo $\mathcal{O}_X$-Módulos localmente libres
de rango finito,

\oldpage[0\textsubscript{I}]{49}

y cuando hablemos de haces localmente libres sin precisar, se
sobreentenderá que son de \emph{rango finito}.
\end{env}

\begin{env}[5.4.2]
\label{0.5.4.2-fr}
Si $\mathcal{L}$, $\mathcal{F}$ son dos $\mathcal{O}_X$-Módulos, se tiene un
homomorfismo canónico funtorial
\[
  \check{\mathcal{L}}\otimes_{\mathcal{O}_X}\mathcal{F}
  =\mathcal{H}\!om_{\mathcal{O}_X}(\mathcal{L},\mathcal{O}_X)
   \otimes_{\mathcal{O}_X}\mathcal{F}
  \longrightarrow
  \mathcal{H}\!om_{\mathcal{O}_X}(\mathcal{L},\mathcal{F})
  \tag{5.4.2.1}\label{0.5.4.2.1-fr}
\]
definido de la manera siguiente: para todo abierto $U$, a todo par $(u,t)$, donde
\[
  u\in\Gamma\bigl(U,
    \mathcal{H}\!om_{\mathcal{O}_X}(\mathcal{L},\mathcal{O}_X)\bigr)
  =\operatorname{Hom}(\mathcal{L}|U,\mathcal{O}_X|U)
  \quad\text{y}\quad
  t\in\Gamma(U,\mathcal{F}),
\]
se hace corresponder el elemento de
$\operatorname{Hom}(\mathcal{L}|U,\mathcal{F}|U)$ que, para todo $x\in U$,
hace corresponder a $s_x\in\mathcal{L}_x$ el elemento $u_x(s_x)t_x$ de
$\mathcal{F}_x$. Si $\mathcal{L}$ es \emph{localmente libre de rango finito},
este homomorfismo es \emph{biyectivo}; como la propiedad es local, se puede,
en efecto, limitar al caso en que $\mathcal{L}=\mathcal{O}_X^n$; como para todo
$\mathcal{O}_X$-Módulo $\mathcal{G}$,
$\mathcal{H}\!om_{\mathcal{O}_X}(\mathcal{O}_X^n,\mathcal{G})$ es
canónicamente isomorfo a $\mathcal{G}^n$, se reduce al caso
$\mathcal{L}=\mathcal{O}_X$, que es inmediato.
\end{env}

\begin{env}[5.4.3]
\EGAReaderTarget{canon-ref.0.5.4.3-fr}{0.5.4.3-fr}
Si $\mathcal{L}$ es invertible, lo mismo ocurre con su dual
$\check{\mathcal{L}}=
\mathcal{H}\!om_{\mathcal{O}_X}(\mathcal{L},\mathcal{O}_X)$, pues se reduce
inmediatamente (siendo la cuestión local) al caso $\mathcal{L}=\mathcal{O}_X$. Además,
se tiene un isomorfismo canónico
\[
  \mathcal{H}\!om_{\mathcal{O}_X}(\mathcal{L},\mathcal{O}_X)
  \otimes_{\mathcal{O}_X}\mathcal{L}
  \xrightarrow{\sim}
  \mathcal{O}_X
  \tag{5.4.3.1}\label{0.5.4.3.1-fr}
\]
en efecto, según (5.4.2), basta definir un isomorfismo canónico
$\mathcal{H}\!om_{\mathcal{O}_X}(\mathcal{L},\mathcal{L})
\xrightarrow{\sim}\mathcal{O}_X$. Ahora bien, para \emph{todo}
$\mathcal{O}_X$-Módulo $\mathcal{F}$, se tiene un homomorfismo canónico
$\mathcal{O}_X\xrightarrow{\sim}
\mathcal{H}\!om_{\mathcal{O}_X}(\mathcal{F},\mathcal{F})$ (5.3.7). Queda
por probar que si $\mathcal{F}=\mathcal{L}$ es invertible, este homomorfismo
es biyectivo, y como la cuestión es local, se reduce al caso
$\mathcal{L}=\mathcal{O}_X$, que es inmediato.

Por lo que precede, se pone
$\mathcal{L}^{-1}=
\mathcal{H}\!om_{\mathcal{O}_X}(\mathcal{L},\mathcal{O}_X)$, y se dice que
$\mathcal{L}^{-1}$ es el \emph{inverso} de $\mathcal{L}$. La terminología de
\og haz invertible\fg{} puede justificarse de la manera siguiente cuando $X$ se
reduce a un punto y $\mathcal{O}_X$ es un anillo \emph{local} $A$ de ideal
maximal $\mathfrak m$; si $M$ y $M'$ son dos $A$-módulos (siendo $M$ de
tipo finito) tales que $M\otimes_A M'$ es isomorfo a $A$, como
$(A/\mathfrak m)\otimes_A(M\otimes_A M')$ se identifica con
$(M/\mathfrak mM)\otimes_{A/\mathfrak m}(M'/\mathfrak mM')$, este último
producto tensorial de espacios vectoriales sobre el cuerpo $A/\mathfrak m$ es
isomorfo a $A/\mathfrak m$, lo que exige que $M/\mathfrak mM$ y
$M'/\mathfrak mM'$ sean de dimensión 1. Para todo elemento $z\in M$
que no pertenece a $\mathfrak mM$, se tiene, pues, $M=Az+\mathfrak mM$, lo que
implica $M=Az$ según el lema de Nakayama, siendo $M$ de tipo finito.
Por lo demás, como el anulador de $z$ anula $M\otimes_A M'$, isomorfo a
$A$, este anulador es $\{0\}$, y $M$ es, por consiguiente, \emph{isomorfo a $A$}.
En el caso general, esto muestra que si $\mathcal{L}$ es un
$\mathcal{O}_X$-Módulo de tipo finito, tal que existe un
$\mathcal{O}_X$-Módulo $\mathcal{F}$ para el cual
$\mathcal{L}\otimes_{\mathcal{O}_X}\mathcal{F}$ es isomorfo a
$\mathcal{O}_X$, y si además los anillos $\mathcal{O}_x$ son anillos
locales, entonces $\mathcal{L}_x$ es un $\mathcal{O}_x$-módulo isomorfo a
$\mathcal{O}_x$ para todo $x\in X$. Si se supone que $\mathcal{O}_X$ y $\mathcal{L}$ son
\emph{coherentes}, se concluye, pues, que $\mathcal{L}$ es invertible
en virtud de (5.2.7).
\end{env}

\begin{env}[5.4.4]
\label{0.5.4.4-fr}
Si $\mathcal{L}$ y $\mathcal{L}'$ son dos $\mathcal{O}_X$-Módulos
invertibles, lo mismo ocurre con
$\mathcal{L}\otimes_{\mathcal{O}_X}\mathcal{L}'$, pues, siendo la cuestión
local, se puede suponer que $\mathcal{L}=\mathcal{O}_X$, y el resultado es
entonces trivial. Para todo entero $n\geq1$, se designa por
$\mathcal{L}^{\otimes n}$ el producto tensorial de $n$ haces

\oldpage[0\textsubscript{I}]{50}

idénticos a $\mathcal{L}$; se pone por convenio
$\mathcal{L}^{\otimes0}=\mathcal{O}_X$ y para $n\geq1$,
$\mathcal{L}^{\otimes(-n)}=(\mathcal{L}^{-1})^{\otimes n}$. Con estas
notaciones, existe entonces un \emph{isomorfismo canónico funtorial}
\[
  \mathcal{L}^{\otimes m}\otimes_{\mathcal{O}_X}\mathcal{L}^{\otimes n}
  \xrightarrow{\sim}
  \mathcal{L}^{\otimes(n+m)}
  \tag{5.4.4.1}\label{0.5.4.4.1-fr}
\]
cualesquiera que sean los enteros racionales $m,n$: en efecto, en virtud de las
definiciones, se reduce inmediatamente al caso en que $m=-1$, $n=1$, y
el isomorfismo en cuestión ha sido entonces definido en (5.4.3).
\end{env}

\begin{env}[5.4.5]
\EGAReaderTarget{canon-ref.0.5.4.5-fr}{0.5.4.5-fr}
Sea $f:Y\to X$ un morfismo de espacios anillados. Si $\mathcal{L}$ es un
$\mathcal{O}_X$-Módulo localmente libre (resp. invertible),
$f^*(\mathcal{L})$ es un $\mathcal{O}_Y$-Módulo localmente libre (resp.
invertible): esto resulta inmediatamente de que las imágenes inversas de dos
$\mathcal{O}_X$-Módulos localmente isomorfos son localmente isomorfas, de
que $f^*$ conmuta con las sumas directas finitas y de que
$f^*(\mathcal{O}_X)=\mathcal{O}_Y$ (4.3.4). Además, se sabe que se tiene un
homomorfismo canónico funtorial
\[
  f^*(\check{\mathcal{L}})
  \longrightarrow
  \bigl(f^*(\mathcal{L})\bigr)\check{}
\]
(4.4.6), y cuando $L$ es localmente libre, este homomorfismo es
\emph{biyectivo}: en efecto, se reduce de nuevo al caso en que
$\mathcal{L}=\mathcal{O}_X$, que es trivial. Se concluye que si
$\mathcal{L}$ es invertible, $f^*(\mathcal{L}^{\otimes n})$ se identifica
canónicamente con $(f^*(\mathcal{L}))^{\otimes n}$ para todo entero racional
$n$.
\end{env}

\begin{env}[5.4.6]
\label{0.5.4.6-fr}
Sea $\mathcal{L}$ un $\mathcal{O}_X$-Módulo invertible; se designa por
$\Gamma_*(X,\mathcal{L})$ o simplemente $\Gamma_*(\mathcal{L})$ el grupo
abeliano suma directa
$\bigoplus_{n\in\mathbb{Z}}\Gamma(X,\mathcal{L}^{\otimes n})$; se lo dota
de una estructura de \emph{anillo graduado}, haciendo corresponder al par
$(s_n,s_m)$, donde $s_n\in\Gamma(X,\mathcal{L}^{\otimes n})$,
$s_m\in\Gamma(X,\mathcal{L}^{\otimes n})$, la sección de
$\mathcal{L}^{\otimes(n+m)}$ sobre $X$ que corresponde canónicamente
(5.4.4.1) a la sección $s_n\otimes s_m$ de
$\mathcal{L}^{\otimes n}\otimes_{\mathcal{O}_X}\mathcal{L}^{\otimes m}$;
la asociatividad de esta multiplicación se verifica de manera inmediata. Es
claro que $\Gamma_*(X,\mathcal{L})$ es un funtor covariante en
$\mathcal{L}$, con valores en la categoría de los anillos graduados.

Si ahora $\mathcal{F}$ es un $\mathcal{O}_X$-Módulo cualquiera, se pone
\[
  \Gamma_*(\mathcal{L},\mathcal{F})
  =\bigoplus_{n\in\mathbb{Z}}
   \Gamma\bigl(X,
     \mathcal{F}\otimes_{\mathcal{O}_X}\mathcal{L}^{\otimes n}\bigr).
\]
Se dota a este grupo abeliano de una estructura de \emph{módulo graduado} sobre
el anillo graduado $\Gamma_*(\mathcal{L})$ de la manera siguiente: al par
$(s_n,u_m)$, donde $s_n\in\Gamma(X,\mathcal{L}^{\otimes n})$ y
$u_m\in\Gamma(X,
\mathcal{F}\otimes\mathcal{L}^{\otimes m})$, se hace
corresponder la sección de
$\mathcal{F}\otimes\mathcal{L}^{\otimes(m+n)}$ que corresponde
canónicamente (5.4.4.1) a $s_n\otimes u_m$; la verificación de los axiomas de
los módulos es inmediata. Para $X$ y $\mathcal{L}$ fijos,
$\Gamma_*(\mathcal{L},\mathcal{F})$ es un funtor covariante en
$\mathcal{F}$ con valores en la categoría de los
$\Gamma_*(\mathcal{L})$-módulos graduados; para $X$ y $\mathcal{F}$ fijos,
es un funtor covariante en $\mathcal{L}$ con valores en la categoría de los
grupos abelianos.

Si $f:Y\to X$ es un morfismo de espacios anillados, el homomorfismo canónico
(4.4.3.2)
\[
  \rho:\mathcal{L}^{\otimes n}
  \longrightarrow
  f_*\bigl(f^*(\mathcal{L}^{\otimes n})\bigr)
\]
define un homomorfismo de grupos abelianos
\[
  \Gamma(X,\mathcal{L}^{\otimes n})
  \longrightarrow
  \Gamma\bigl(Y,f^*(\mathcal{L}^{\otimes n})\bigr),
\]
y como
$f^*(\mathcal{L}^{\otimes n})=(f^*(\mathcal{L}))^{\otimes n}$, resulta de las
definiciones de los homomorfismos canónicos (4.4.3.2) y (5.4.4.1) que los
homomorfismos anteriores definen un \emph{homomorfismo funtorial
de anillos graduados}
$\Gamma_*(\mathcal{L})\to\Gamma_*(f^*(\mathcal{L}))$. El mismo homomorfismo
canónico (4.4.3) define asimismo un homomorfismo de grupos abelianos
\[
  \Gamma\bigl(X,
    \mathcal{F}\otimes\mathcal{L}^{\otimes n}\bigr)
  \longrightarrow
  \Gamma\bigl(Y,
    f^*(\mathcal{F}\otimes\mathcal{L}^{\otimes n})\bigr),
\]
y como
\[
  f^*(\mathcal{F}\otimes\mathcal{L}^{\otimes n})
  =f^*(\mathcal{F})\otimes
   (f^*(\mathcal{L}))^{\otimes n}
  \quad\text{(4.3.3.1),}
\]

\oldpage[0\textsubscript{I}]{51}

estos homomorfismos (para $n$ variable) definen un
\emph{dihomomorfismo de módulos graduados}
\[
  \Gamma_*(\mathcal{L},\mathcal{F})
  \longrightarrow
  \Gamma_*(f^*(\mathcal{L}),f^*(\mathcal{F})).
\]
\end{env}

\begin{env}[5.4.7]
\EGAReaderTarget{canon-ref.0.5.4.7-fr}{0.5.4.7-fr}
Se puede mostrar que existe un \emph{conjunto} $\mathfrak{M}$ (denotado también
por $\mathfrak{M}(X)$) de $\mathcal{O}_X$-Módulos invertibles tal que todo
$\mathcal{O}_X$-Módulo invertible es isomorfo a un elemento de
$\mathfrak{M}$ y a uno solo\footnote{Véase el libro en preparación citado en
la Introducción.}; se define en $\mathfrak{M}$ una ley de composición
haciendo corresponder a dos elementos $\mathcal{L}$, $\mathcal{L}'$ de
$\mathfrak{M}$ el único elemento de $\mathfrak{M}$ isomorfo a
$\mathcal{L}\otimes\mathcal{L}'$. Con esta ley de
composición, $\mathfrak{M}$ es un \emph{grupo isomorfo al grupo de
cohomología} $H^1(X,\mathcal{O}_X^*)$, donde $\mathcal{O}_X^*$ es el
subhaz de $\mathcal{O}_X$ tal que $\Gamma(U,\mathcal{O}_X^*)$ es el
grupo de los elementos invertibles del anillo $\Gamma(U,\mathcal{O}_X)$ para
todo abierto $U\subset X$ ($\mathcal{O}_X^*$ es, por tanto, un haz de grupos
abelianos \emph{multiplicativos}).

Para ello se observará que, para todo abierto $U\subset X$, el grupo de las secciones
$\Gamma(U,\mathcal{O}_X^*)$ se identifica canónicamente con el \emph{grupo de los
automorfismos} del $(\mathcal{O}_X|U)$-Módulo $\mathcal{O}_X|U$,
haciendo corresponder la identificación a una sección $\varepsilon$ de
$\mathcal{O}_X^*$ sobre $U$ el automorfismo $u$ de $\mathcal{O}_X|U$
tal que $u_x(s_x)=\varepsilon_xs_x$ para todo $x\in X$ y todo
$s_x\in\mathcal{O}_x$. Sea entonces $\mathfrak{U}=(U_\lambda)$ un recubrimiento
abierto de $X$; el dato, para todo par de índices $(\lambda,\mu)$, de un
automorfismo $\theta_{\lambda\mu}$ de
$\mathcal{O}_X|(U_\lambda\cap U_\mu)$ equivale a darse una
$1$-\emph{cocadena} del recubrimiento $\mathfrak{U}$, con valores en
$\mathcal{O}_X^*$, y decir que los $\theta_{\lambda\mu}$ verifican la
condición de pegado (3.3.1) significa que la cocadena correspondiente es
un \emph{cociclo}. Del mismo modo, el dato, para todo $\lambda$, de un automorfismo
$\omega_\lambda$ de $\mathcal{O}_X|U_\lambda$ equivale al dato de una
$0$-cocadena del recubrimiento $\mathfrak{U}$, con valores en
$\mathcal{O}_X^*$, y su \emph{coborde} corresponde a la familia de los
automorfismos
$(\omega_\lambda|U_\lambda\cap U_\mu)\circ
 (\omega_\mu|U_\lambda\cap U_\mu)^{-1}$.
Se puede hacer corresponder a todo $1$-cociclo de $\mathfrak{U}$ con valores
en $\mathcal{O}_X^*$ el elemento de $\mathfrak{M}$ isomorfo al
$\mathcal{O}_X$-Módulo invertible obtenido por pegado a partir de la
familia de automorfismos $(\theta_{\lambda\mu})$ correspondiente a este cociclo,
y a dos cociclos cohomólogos corresponden dos elementos iguales de
$\mathfrak{M}$ (3.3.2); en otras palabras, se ha definido así una aplicación
$\varphi_{\mathfrak{U}}:H^1(\mathfrak{U},\mathcal{O}_X^*)\to\mathfrak{M}$.
Además, si $\mathfrak{B}$ es un segundo recubrimiento abierto de $X$, más fino
que $\mathfrak{U}$, el diagrama
\[
\begin{tikzcd}[column sep=huge,row sep=large]
H^1(\mathfrak{U},\mathcal{O}_X^*)
  \arrow[dd]
  \arrow[rrd,"\varphi_{\mathfrak{U}}"] & & \\
& & \mathfrak{M} \\
H^1(\mathfrak{B},\mathcal{O}_X^*)
  \arrow[rru,swap,"\varphi_{\mathfrak{B}}"] & &
\end{tikzcd}
\]
donde la flecha vertical es el homomorfismo canónico (G, II, 5.7), es
conmutativo, como resulta de (3.3.3). Por paso al límite inductivo, se
obtiene, pues, una aplicación
$H^1(X,\mathcal{O}_X^*)\to\mathfrak{M}$, identificándose el grupo de cohomología de \v{C}ech
$\check{H}^1(X,\mathcal{O}_X^*)$, como se sabe, con el primer grupo
de cohomología $H^1(X,\mathcal{O}_X^*)$ (G, II, 5.9, cor. del teorema~5.9.1).
Esta aplicación es \emph{sobreyectiva}: en efecto, por definición, para todo
$\mathcal{O}_X$-Módulo invertible $\mathcal{L}$, hay un recubrimiento abierto
$(U_\lambda)$ de $X$ tal que $\mathcal{L}$ se obtiene por pegado de los
haces $\mathcal{O}_X|U_\lambda$ (3.3.1). Es también
\emph{inyectiva}, pues basta probarlo para las aplicaciones
$H^1(\mathfrak{U},\mathcal{O}_X)\to\mathfrak{M}$, y esto resulta entonces de
(3.3.2). Queda por mostrar que

\oldpage[0\textsubscript{I}]{52}

la biyección así definida es un homomorfismo de grupos. Ahora bien, dados
dos $\mathcal{O}_X$-Módulos invertibles $\mathcal{L}$, $\mathcal{L}'$, hay
un recubrimiento abierto $(U_\lambda)$ tal que $\mathcal{L}|U_\lambda$ y
$\mathcal{L}'|U_\lambda$ son isomorfos a $\mathcal{O}_X|U_\lambda$ para
todo $\lambda$; hay, pues, para cada índice $\lambda$ un elemento
$a_\lambda$ (resp. $a'_\lambda$) de $\Gamma(U_\lambda,\mathcal{L})$ (resp.
$\Gamma(U_\lambda,\mathcal{L}'))$ tal que los elementos de
$\Gamma(U_\lambda,\mathcal{L})$ (resp. $\Gamma(U_\lambda,\mathcal{L}'))$
son los $s_\lambda\cdot a_\lambda$ (resp.
$s_\lambda\cdot a'_\lambda$), donde $s_\lambda$ recorre
$\Gamma(U_\lambda,\mathcal{O}_X)$. Los cociclos correspondientes
$(\varepsilon_{\lambda\mu})$, $(\varepsilon'_{\lambda\mu})$ son tales que
$s_\lambda\cdot a_\lambda=s_\mu\cdot a_\mu$ (resp.
$s_\lambda\cdot a'_\lambda=s_\mu\cdot a'_\mu$) sobre
$U_\lambda\cap U_\mu$ es equivalente a
$s_\lambda=\varepsilon_{\lambda\mu}s_\mu$ (resp.
$s_\lambda=\varepsilon'_{\lambda\mu}s_\mu$) sobre
$U_\lambda\cap U_\mu$. Como las secciones de
$\mathcal{L}\otimes_{\mathcal{O}_X}\mathcal{L}'$ sobre $U_\lambda$
son las sumas finitas de los
$s_\lambda s'_\lambda\cdot(a_\lambda\otimes a'_\lambda)$ donde $s_\lambda$ y
$s'_\lambda$ recorren $\Gamma(U_\lambda,\mathcal{O}_X)$, es claro que
el cociclo $(\varepsilon_{\lambda\mu}\varepsilon'_{\lambda\mu})$ corresponde
a $\mathcal{L}\otimes_{\mathcal{O}_X}\mathcal{L}'$, lo que termina la
demostración\footnote{Para una forma general de este resultado, véase el libro
citado en la nota de la p.~51.}.
\end{env}

\begin{env}[5.4.8]
\label{0.5.4.8-fr}
Sea $f=(\psi,\omega)$ un morfismo $Y\to X$ de espacios anillados. El funtor
$f^*(\mathcal{L})$ en la categoría de los $\mathcal{O}_X$-Módulos libres
define una aplicación (denotada también por $f^*$ por abuso de lenguaje) del conjunto
$\mathfrak{M}(X)$ en el conjunto $\mathfrak{M}(Y)$. Por otra parte, se tiene un
homomorfismo canónico (T, 3.2.2)
\[
  H^1(X,\mathcal{O}_X^*)\longrightarrow H^1(Y,\mathcal{O}_Y^*).
  \tag{5.4.8.1}\label{0.5.4.8.1-fr}
\]
Cuando se identifican canónicamente (5.4.7) $\mathfrak{M}(X)$ y
$H^1(X,\mathcal{O}_X^*)$ (resp. $\mathfrak{M}(Y)$ y
$H^1(Y,\mathcal{O}_Y^*)$), el homomorfismo (5.4.8.1) \emph{se identifica con
la aplicación $f^*$}. En efecto, si $\mathcal{L}$ procede de un cociclo
$(\varepsilon_{\lambda\mu})$ correspondiente a un recubrimiento abierto
$(U_\lambda)$ de $X$, basta mostrar que $f^*(\mathcal{L})$ procede
de un cociclo cuya clase de cohomología es imagen por (5.4.8.1) de la de
$(\varepsilon_{\lambda\mu})$. Ahora bien, si $\theta_{\lambda\mu}$ es
el automorfismo de $\mathcal{O}_X|(U_\lambda\cap U_\mu)$ que corresponde a
$\varepsilon_{\lambda\mu}$, es claro que $f^*(\mathcal{L})$ se obtiene por
pegado de los $\mathcal{O}_Y|\psi^{-1}(U_\lambda)$ mediante los
automorfismos $f^*(\theta_{\lambda\mu})$, y basta, pues, verificar que
estos últimos corresponden al cociclo
$(\omega^\sharp(\varepsilon_{\lambda\mu}))$, lo que resulta inmediatamente de las
definiciones (se ha identificado $\varepsilon_{\lambda\mu}$ con su imagen canónica
por $\rho$ (3.7.2), sección de $\psi^*(\mathcal{O}_X^*)$ sobre
$\psi^{-1}(U_\lambda\cap U_\mu)$).
\end{env}

\begin{env}[5.4.9]
\label{0.5.4.9-fr}
Sean $\mathcal{E}$, $\mathcal{F}$ dos $\mathcal{O}_X$-Módulos,
suponiéndose $\mathcal{F}$ \emph{localmente libre}, y sea $\mathcal{G}$ un
$\mathcal{O}_X$-Módulo \emph{extensión de $\mathcal{F}$ por $\mathcal{E}$},
es decir, tal que existe una sucesión exacta
\[
  0\longrightarrow\mathcal{E}\xrightarrow{j}\mathcal{G}
  \xrightarrow{p}\mathcal{F}\longrightarrow 0.
\]
Entonces, para todo $x\in X$, existe una vecindad abierta $U$ de $x$ tal que
$\mathcal{G}|U$ es isomorfo a la \emph{suma directa}
$\mathcal{E}|U\oplus\mathcal{F}|U$. Se puede, en efecto, limitar al caso en que
$\mathcal{F}=\mathcal{O}_X^n$; sean $e_i$ ($1\leq i\leq n$) las secciones
canónicas (5.5.5) de $\mathcal{O}_X^n$; existe entonces una vecindad abierta
$U$ de $x$ y $n$ secciones $s_i$ de $\mathcal{G}$ sobre $U$ tales que
$p(s_i|U)=e_i|U$ para $1\leq i\leq n$. Dicho esto, sea $f$
el homomorfismo $\mathcal{F}|U\to\mathcal{G}|U$ definido por las secciones
$s_i|U$ (5.1.1). Es inmediato que para todo abierto $V\subset U$, y toda
sección $s\in\Gamma(V,\mathcal{G})$ se tiene
$s-f(p(s))\in\Gamma(V,\mathcal{E})$, de donde nuestra afirmación.
\end{env}

\begin{env}[5.4.10]
\label{0.5.4.10-fr}
Sean $f:X\to Y$ un morfismo de espacios anillados, $\mathcal{F}$ un
$\mathcal{O}_X$-Módulo, $\mathcal{L}$ un $\mathcal{O}_Y$-Módulo localmente
libre de rango finito. Existe entonces un isomorfismo canónico
\[
  f_*(\mathcal{F})\otimes_{\mathcal{O}_Y}\mathcal{L}
  \xrightarrow{\sim}
  f_*\bigl(\mathcal{F}\otimes_{\mathcal{O}_X}f^*(\mathcal{L})\bigr)
  \tag{5.4.10.1}\label{0.5.4.10.1-fr}
\]

\oldpage[0\textsubscript{I}]{53}

En efecto, para todo $\mathcal{O}_Y$-Módulo $\mathcal{L}$, se tiene un
homomorfismo canónico
\[
  f_*(\mathcal{F})\otimes_{\mathcal{O}_Y}\mathcal{L}
  \xrightarrow{1\otimes\rho}
  f_*(\mathcal{F})\otimes_{\mathcal{O}_Y}f_*(f^*(\mathcal{L}))
  \xrightarrow{\alpha}
  f_*\bigl(\mathcal{F}\otimes_{\mathcal{O}_X}f^*(\mathcal{L})\bigr)
\]
siendo $\rho$ el homomorfismo (4.4.3.2) y $\alpha$ el homomorfismo (4.2.2.1).
Para mostrar que cuando $\mathcal{L}$ es localmente libre este homomorfismo
es biyectivo, basta, siendo la cuestión local sobre $Y$, considerar el
caso en que $\mathcal{L}=\mathcal{O}_Y^n$; además, como $f_*$ y $f^*$ conmutan con las
sumas directas finitas, se puede suponer $n=1$, y en este caso la proposición
resulta inmediatamente de las definiciones y de la relación
$f^*(\mathcal{O}_Y)=\mathcal{O}_X$.
\end{env}

